% ============================================================
% 软件开发实践：沟通讨论记录 LaTeX 模板
% 编译方式：XeLaTeX
% 用途：将“成员 ↔ GPT”的关键讨论整理成可读、可复用的对话记录
% ============================================================

\documentclass[12pt,a4paper]{ctexart}

\usepackage[
  top=2.5cm,
  bottom=2.5cm,
  left=3cm,
  right=2.5cm
]{geometry}

\usepackage{xcolor}
\usepackage[most]{tcolorbox}
\usepackage{fontawesome5}
\usepackage{hyperref}
\usepackage{enumitem}

\hypersetup{
  colorlinks=true,
  linkcolor=blue!65!black,
  urlcolor=blue!65!black
}

% -------------------- 配色 --------------------
\definecolor{memberbg}{RGB}{239,246,255}
\definecolor{memberborder}{RGB}{96,165,250}
\definecolor{gptbg}{RGB}{240,253,250}
\definecolor{gptborder}{RGB}{45,150,120}
\definecolor{decisionbg}{RGB}{248,250,252}
\definecolor{decisionborder}{RGB}{148,163,184}

% -------------------- 对话框 --------------------
% 成员消息：左侧
\newtcolorbox{membermsg}[1]{
  enhanced,
  breakable,
  colback=memberbg,
  colframe=memberborder,
  boxrule=0.6pt,
  arc=2.5mm,
  left=4mm,
  right=4mm,
  top=2.5mm,
  bottom=2.5mm,
  width=0.86\linewidth,
  right skip=0.14\linewidth,
  before skip=3mm,
  after skip=2mm,
  title={\faUserCircle\quad #1},
  fonttitle=\small\bfseries,
  coltitle=black,
  colbacktitle=memberbg,
  boxed title style={
    boxrule=0pt,
    colframe=memberbg,
    colback=memberbg
  },
  attach boxed title to top left={xshift=2mm,yshift=-2mm}
}

% GPT 消息：右侧
\newtcolorbox{gptmsg}{
  enhanced,
  breakable,
  colback=gptbg,
  colframe=gptborder,
  boxrule=0.6pt,
  arc=2.5mm,
  left=4mm,
  right=4mm,
  top=2.5mm,
  bottom=2.5mm,
  width=0.86\linewidth,
  left skip=0.14\linewidth,
  before skip=3mm,
  after skip=2mm,
  title={\faRobot\quad GPT},
  fonttitle=\small\bfseries,
  coltitle=black,
  colbacktitle=gptbg,
  boxed title style={
    boxrule=0pt,
    colframe=gptbg,
    colback=gptbg
  },
  attach boxed title to top left={xshift=2mm,yshift=-2mm}
}

% 讨论结论
\newtcolorbox{decisionbox}{
  enhanced,
  breakable,
  colback=decisionbg,
  colframe=decisionborder,
  boxrule=0.6pt,
  arc=1.5mm,
  left=4mm,
  right=4mm,
  top=3mm,
  bottom=3mm,
  before skip=3mm,
  after skip=5mm
}

% -------------------- 文档信息 --------------------
\title{软件开发实践\\沟通讨论记录样例}
\author{WEMOVE 项目组}
\date{2026 年 9 月}

\begin{document}

\maketitle

\section{记录说明}

本模板用于保留软件开发实践中的关键沟通过程。每个主题保留能够体现“提出问题、分析方案、形成结论”的核心对话，并在对话末尾给出明确的讨论结论。无需完整复制全部聊天内容。

\section{成员 D 与 GPT 的项目沟通样例}

\subsection{模块是否需要独立应用}

\noindent
\textbf{讨论时间：}2026 年 9 月 5 日至 9 月 8 日

\begin{membermsg}{成员 D}
经销合作模块是否需要单独创建一套后端应用、路由、数据库或登录体系？
\end{membermsg}

\begin{gptmsg}
当前项目已经具备统一的后端、身份与数据基础设施。更合适的做法是在现有 Spring Boot 工程中扩展 \texttt{wemove.dealership} 模块，前端接入现有 Vue 应用的统一路由，并继续复用身份认证、幂等、限频和数据库迁移能力。
\end{gptmsg}

\begin{membermsg}{成员 D}
商品、文件和运营统计这些能力也都直接接现有模块吗？
\end{membermsg}

\begin{gptmsg}
是。经销合作模块应调用已有能力：商品信息沿用商品模块的事实数据，受保护资料沿用文件模块，审计与待办统计沿用运营模块。这样可以避免同一份业务事实被多套实现重复维护。
\end{gptmsg}

\begin{decisionbox}
\textbf{讨论结论：}
经销合作模块继续作为统一系统中的业务模块实现，不新增第二套用户、登录、商品、文件或审计体系。
\end{decisionbox}


\subsection{本地登录出现 ORIGIN\_DENIED}

\begin{membermsg}{成员 D}
本地启动前后端后，登录接口返回 \texttt{ORIGIN\_DENIED}，这个问题一般从哪里排查？
\end{membermsg}

\begin{gptmsg}
优先核对请求的 Origin。Origin 由协议、主机和端口共同组成，因此 \texttt{http://localhost:5173} 与 \texttt{http://127.0.0.1:5173} 会被视为两个不同来源。
\end{gptmsg}

\begin{membermsg}{成员 D}
如果前端现在使用的是 \texttt{127.0.0.1:5173} 呢？
\end{membermsg}

\begin{gptmsg}
可以把该来源加入 \texttt{ALLOWED\_ORIGINS}，然后重启后端使配置重新加载。团队开发时最好统一使用 \texttt{http://localhost:5173}，可以减少环境差异带来的问题。
\end{gptmsg}

\begin{decisionbox}
\textbf{讨论结论：}
开发环境统一使用 \texttt{http://localhost:5173}；需要使用其他来源时，将完整 Origin 加入允许列表，并在修改配置后重启后端。
\end{decisionbox}


\subsection{数据库字段与敏感信息边界}

\begin{membermsg}{成员 D}
排查用户账户状态时，我应该查询 \texttt{status} 字段吗？
\end{membermsg}

\begin{gptmsg}
当前 \texttt{users} 表的账户状态字段是 \texttt{account\_status}。排查时只选择实际需要的字段，例如邮箱、角色和账户状态。
\end{gptmsg}

\begin{membermsg}{成员 D}
那 \texttt{password\_hash} 可以一起查出来方便确认账号吗？
\end{membermsg}

\begin{gptmsg}
不需要。密码哈希与账户是否启用无关，也不应进入截图、聊天记录、日志或问题文档。查询和记录时都应控制在排查所需的最小字段范围内。
\end{gptmsg}

\begin{decisionbox}
\textbf{讨论结论：}
账户状态使用 \texttt{account\_status} 判断；沟通记录和排查材料中不得记录 \texttt{password\_hash} 等敏感字段。
\end{decisionbox}


\section{后续填写方式}

复制一个讨论主题时，只需要保留以下结构：

\begin{enumerate}[label=\arabic*.]
  \item 使用 \verb|\subsection{讨论主题}| 标记本次讨论的核心问题；
  \item 用 \verb|membermsg| 写成员提出的问题或方案；
  \item 用 \verb|gptmsg| 写 GPT 的关键分析与建议；
  \item 用 \verb|decisionbox| 写最终采用的方案、限制或待办事项。
\end{enumerate}

建议每个主题保留 2--6 轮关键消息。较长的原始聊天可以先提炼，再进入正式沟通记录。

\end{document}
